\chapter{REST API：Controller、参数校验、统一错误处理}

最后一关：把数据库里的数据通过 REST API 提供出去。做完之后，\code{src/main/resources/static/map.html} 这个地图页面（已经写好了）就能显示所有站点。

\partinfo{%
  \code{station/StationQueryService.java}：6.1\\
  \code{station/web/StationController.java}：6.2、6.3\\
  \code{web/GlobalExceptionHandler.java}：6.3%
}{\code{./mvnw test -Dtest='Part6*'}，最后 \code{./mvnw test -Dtest=FinalApplicationTest}}

\section{知识点}

\subsection{分层}

\begin{plaincode}
HTTP 请求 --> Controller --> Service --> Repository --> 数据库
             (处理 HTTP)    (业务逻辑)   (数据访问)
\end{plaincode}

\begin{itemize}
  \item \textbf{Controller} 只负责 HTTP 相关的事：解析参数、校验、把结果转成 JSON、设置状态码。不写业务逻辑。
  \item \textbf{Service} 写业务逻辑，组合多个 Repository 的数据。不知道 HTTP 的存在。
  \item \textbf{Repository} 只负责读写数据库。
\end{itemize}

这一关有三种“站点”对象，各自负责一层：

\begin{center}
\begin{tabular}{lll}
\toprule
类 & 层 & 作用 \\
\midrule
\code{StationEntity} & 数据库 & 对应表结构 \\
\code{Station} & 业务 & 业务代码使用的对象（Part 1） \\
\code{StationResponse} & API & 返回给前端的 JSON 结构（已经写好了） \\
\bottomrule
\end{tabular}
\end{center}

\textbf{为什么不直接把 Entity 返回给前端？}表结构一改，API 就跟着变，前端会出问题；而且 Entity 里可能有不该暴露的字段。

\subsection{避免 N+1 查询}

\code{listStations()} 要返回每个站点和它最新的可用量。最直接的写法是：

\begin{java}
for (StationEntity s : stationRepository.findAll()) {
    availabilityRepository.findLatest(s.getNumber());   // 每个站点查一次
}
\end{java}

113 个站点就是 1 + 113 = 114 条 SQL，这就是 \textbf{N+1 问题}。更好的做法：用 Part 4.5 的查询，一次拿到所有站点的最新数据，放进一个 \code{Map<站点编号, AvailabilityEntity>}，再逐个查 Map。这样总共只有 2 条 SQL。

\begin{java}
Map<Long, BookEntity> byIsbn = books.stream()
        .collect(Collectors.toMap(b -> b.getIsbn(), b -> b));  // key 是 isbn，value 是对象本身
byIsbn.get(isbn);                                             // 找不到时返回 null
\end{java}

\code{b -> b} 也可以写成 \code{Function.identity()}。按字段排序查全部：\code{stationRepository.findAll(Sort.by("number"))}。

\subsection{Controller}

\begin{java}
@RestController                          // 返回值自动转成 JSON
@RequestMapping("/api/v1/books")         // 这个类里所有接口的公共路径
public class BookController {

    @GetMapping                          // GET /api/v1/books
    public List<BookResponse> list() { ... }

    @GetMapping("/{isbn}")               // GET /api/v1/books/123
    public BookResponse get(@PathVariable long isbn) { ... }

    @GetMapping("/search")               // GET /api/v1/books/search?q=dune&limit=5
    public List<BookResponse> search(
            @RequestParam String q,                                // 必填，缺了返回 400
            @RequestParam(defaultValue = "10") int limit) { ... }  // 可选，有默认值
}
\end{java}

record 转成 JSON 时，字段名就是 record 的字段名（camelCase）：\code{StationResponse} 的 \code{availableBikes} 会变成 \code{"availableBikes": 3}。

\subsection{参数校验}

在参数上加 \code{jakarta.validation.constraints} 包里的注解，Spring 会在调用方法之前检查，不满足就直接返回 400，方法根本不会执行：

\begin{java}
public List<BookResponse> search(
        @RequestParam @Size(min = 2) String q,
        @RequestParam(defaultValue = "10") @Min(1) @Max(50) int limit) { ... }
\end{java}

\begin{center}
\begin{tabular}{ll}
\toprule
注解 & 用在 \\
\midrule
\code{@Min} / \code{@Max} & 整数 \\
\code{@DecimalMin("-90")} / \code{@DecimalMax("90")} & 小数（值写成字符串） \\
\code{@NotBlank}、\code{@Size}、\code{@Pattern} & 字符串 \\
\bottomrule
\end{tabular}
\end{center}

\subsection{统一错误处理和 ProblemDetail}

Controller 里抛出的异常，可以在一个全局的地方统一转成 HTTP 响应：

\begin{java}
@RestControllerAdvice                                  // 对所有 Controller 生效
public class GlobalExceptionHandler extends ResponseEntityExceptionHandler {

    @ExceptionHandler(BookNotFoundException.class)     // 处理这种异常
    public ProblemDetail handleNotFound(BookNotFoundException e) {
        ProblemDetail problem =
                ProblemDetail.forStatusAndDetail(HttpStatus.NOT_FOUND, e.getMessage());
        problem.setTitle("Book not found");
        problem.setProperty("isbn", e.getIsbn());      // 可以加自定义字段
        return problem;
    }
}
\end{java}

\code{ProblemDetail} 是 RFC 9457 定义的标准错误格式，返回的 Content-Type 是 \code{application/problem+json}：

\begin{jsoncode}
{
  "title": "Book not found",
  "status": 404,
  "detail": "Book 123 does not exist",
  "instance": "/api/v1/books/123",
  "isbn": 123
}
\end{jsoncode}

继承 \code{ResponseEntityExceptionHandler} 的作用：Spring 自己抛的异常（缺少必填参数、参数校验失败、参数类型不对等）也会自动转成 400 的 ProblemDetail，不用一个个自己写。

\subsection{测试是怎么验证的}

\begin{itemize}
  \item \code{Part6QueryServiceTest}：\code{@DataJpaTest} 加真实的内存数据库，测 Service。
  \item \code{Part6ControllerTest}：\code{@WebMvcTest(StationController.class)} 只启动 Web 层，不连数据库。\code{@MockitoBean} 把 \code{StationQueryService} 换成 Mockito 的假对象；\code{MockMvc} 模拟发 HTTP 请求，不用真的启动服务器；\code{jsonPath("$[0].availableBikes")} 检查 JSON 里的值。
  \item \code{FinalApplicationTest}：\code{@SpringBootTest} 启动整个应用，检查所有 Bean 都能正确创建、连在一起。前面任何一个类忘了加 \code{@Service}、\code{@Component}，这个测试都会失败。
\end{itemize}

\section{题目}

\begin{task}{6.1\quad StationQueryService}
\code{StationQueryService} 是一个 Service，类里所有方法都在\textbf{只读事务}里执行（类上加一次就行）。

\textbf{a) \code{listStations()}}：返回所有站点，按 \code{number} 升序。每个站点配上它最新的可用量，用已经写好的 \code{toStation(entity, availability)} 转换（\code{availability} 可以是 null）。\textbf{只能用 2 次查询}：一次查所有站点，一次用 \code{findLatestForAllStations()} 查所有最新可用量。

\begin{example}
\begin{plaincode}
站点：42、33、5
可用量：42 有 08:00（1 辆）和 08:05（3 辆）；33 有 08:05（0 辆）；5 没有记录

listStations()  ->  [5（0 辆, 0 空位）, 33（0 辆）, 42（3 辆, 27 空位）]
\end{plaincode}
\end{example}

\textbf{b) \code{getStation(number)}}：返回单个站点，配上 \code{findLatest(number)} 的结果（可能没有）。站点不存在时抛 \code{StationNotFoundException}（已经写好了）。

\textbf{c) \code{findNearest(lat, lng, limit, requireBikes)}}：用 \code{listStations()} 拿到所有站点，再交给 Part 1 的 \code{StationFinder.findNearest}（\code{requireStands} 传 false）。
\end{task}

\begin{task}{6.2\quad StationController}
\code{StationController} 是一个 REST Controller，公共路径 \code{/api/v1/stations}。

\textbf{a) \code{GET /api/v1/stations}}：返回所有站点，每个用 \code{StationResponse.from(...)} 转换。

\begin{jsoncode}
[
  {
    "number": 42, "name": "SMITHFIELD NORTH", "address": "Smithfield North",
    "lat": 53.349562, "lng": -6.278198, "status": "OPEN",
    "availableBikes": 3, "availableStands": 27, "totalStands": 30
  }
]
\end{jsoncode}

\textbf{b) \code{GET /api/v1/stations/{number}}}：返回单个站点，\code{number} 从路径里取。

\textbf{c) \code{GET /api/v1/stations/nearest}}：

\medskip
\begin{tabular}{lll}
\toprule
参数 & 必填 & 默认值 \\
\midrule
\code{lat} & 是 & \\
\code{lng} & 是 & \\
\code{limit} & 否 & 3 \\
\code{requireBikes} & 否 & false \\
\bottomrule
\end{tabular}

\medskip
每个结果用 \code{NearbyStationResponse.from(...)} 转换。

\begin{example}
\begin{plaincode}
GET /api/v1/stations/nearest?lat=53.3498&lng=-6.2603
-> [{"number": 42, "name": "SMITHFIELD NORTH", ..., "distanceMeters": 1188}]
\end{plaincode}
\end{example}
\end{task}

\begin{task}{6.3\quad 错误处理和参数校验}
\textbf{a) 站点不存在}：完成 \code{GlobalExceptionHandler}，\code{StationNotFoundException} 转成 404 ProblemDetail：

\begin{plaincode}
GET /api/v1/stations/999

404, Content-Type: application/problem+json
{
  "title": "Station not found",
  "status": 404,
  "detail": "Station 999 does not exist",
  "stationNumber": 999,
  ...
}
\end{plaincode}

\textbf{b) 参数校验}：\code{lat} 在 -90 到 90 之间，\code{lng} 在 -180 到 180 之间，\code{limit} 在 1 到 20 之间。不满足时返回 400 ProblemDetail，而且\textbf{不能调用} Service。

\begin{example}
\begin{plaincode}
GET /api/v1/stations/nearest?lat=100&lng=-6.26           ->  400
GET /api/v1/stations/nearest?lat=53.3&lng=-6.2&limit=50  ->  400
GET /api/v1/stations/nearest?lng=-6.26                   ->  400（缺少 lat）
\end{plaincode}
\end{example}
\end{task}

\section{全部完成}

\begin{shell}
./mvnw test                                  # 58 个测试全绿
SCRAPER_ENABLED=true ./mvnw spring-boot:run
\end{shell}

几秒钟后打开 \url{http://localhost:8080/map.html}，地图上应该有 113 个站点。在另一个终端试试这些请求：

\begin{shell}
curl localhost:8080/api/v1/stations/42
curl "localhost:8080/api/v1/stations/nearest?lat=53.3498&lng=-6.2603&requireBikes=true"
curl localhost:8080/api/v1/stations/99999
curl "localhost:8080/api/v1/stations/nearest?lat=200&lng=0"
\end{shell}

\begin{think}
\begin{enumerate}
  \item Controller、Service、Repository 为什么要分层？如果把查询直接写在 Controller 里，会有什么问题？
  \item \code{@Transactional(readOnly = true)} 有什么好处？
  \item \code{@WebMvcTest} 和 \code{@SpringBootTest} 有什么区别？各自适合测什么？
  \item 为什么 API 路径里要带版本号 \code{/v1}？要改成不兼容的格式时怎么办？
  \item 现在每次请求 \code{/api/v1/stations} 都会查两次数据库，但数据 5 分钟才更新一次。怎么优化？（提示：缓存，下一阶段的内容）
\end{enumerate}
\end{think}
